\section{Introduction}
\label{sec:intro}

Multi-task benchmarks are central infrastructure for comparing machine-learning models: a leaderboard has several sub-datasets, each model has a score on each, and an aggregation rule such as the mean or the median produces the ranking. Rankings influence adoption, publication and funding, so a developer has an incentive to push a particular model up. One lever is to fold data related to some sub-datasets into that model's training, which raises its scores there \citep{dominguezolmedo2024ttt,zhou2023cheater}. How much such targeted investment does first place require, and how does the answer depend on the aggregation rule?

We answer this on observed scores, without training any model. A target model may raise its score on any task up to a per-task cap, standing in for what training on related data could buy; all other models are fixed. For each rule we compute the \emph{exact} minimum investment that makes the target strictly first, and report three measures: the total points raised ($L_1$), the number of tasks touched ($K$), and the largest single-task raise ($L_\infty$, the cost of a uniform ``general improvement''). The caps are the modelling assumption on which every conclusion is conditional, so we vary them: a ceiling, the best existing score per task, a fixed gain per task, and a cap calibrated on observed fine-tune-minus-base differences. We also compute the \emph{minimum cap} $h^\ast$ at which first place becomes feasible.

\paragraph{Cost model.} Investment is measured in score points with a linear, task-symmetric cost: a point costs the same on every task and at every level. This is a deliberate simplification, not a claim about what a developer spends; the three measures bracket it ($K$ counts tasks, i.e.\ the fixed cost of a data campaign; $L_\infty$ the cost of a uniform improvement), and Section~\ref{sec:limits} states what it leaves out.

\paragraph{Contributions.} We separate what follows from the definitions from what we measured, and, among the measurements, what holds on both datasets from what depends on the cap, the sample or the data slice. Each claim is followed by its verdict, including those that contradicted what we expected before running the grids.
\begin{itemize}
\item \textbf{Derived (Section~\ref{sec:method}); not empirical.} For the mean, $L_1^\ast=n\,L_\infty^\ast$ whenever no cap binds, so targeting saves no points, only concentrates them in fewer tasks; for the minimum and the median we give closed forms, with $K^\ast\le(n+1)/2$ for the median. Frontier caps make first place reachable under every monotone rule by construction. The solvers agree with the formulas (Section~\ref{sec:exp-theory}).
\item \textbf{Measured, on both datasets.} Contrary to our expectation, the minimum is cheaper than the mean in points for every target under the linear cost (for \cvxHelmQuadraticShare\% and \cvxOllbQuadraticShare\% of targets under a quadratic cost; Appendix~\ref{app:sens}), and relatively cheapest for balanced targets (the cost ratio rises with the coefficient of variation of the target's scores, Spearman $\rho=$~\ctwoHelmRho{} on HELM, \ctwoOllbRho{} on the Open LLM Leaderboard; other dispersion statistics are not stable). The median is not cheaper than the minimum in tasks touched at any task count we tried, again contradicting our expectation (Section~\ref{sec:exp-theory}). Cost grows with rank, and rank alone already explains most of the variance; the gap to the best rival adds explanatory power for the geometric and harmonic means and the median on both datasets and for Borda on HELM, but not for the minimum, the Condorcet requirement, Borda on the Open LLM Leaderboard or plurality on HELM (Section~\ref{sec:exp-rank}); for the first three the gap is close to the cost by construction, so this is a weak test.
\item \textbf{Measured, but conditional on the cap, the sample or the slice.} Whether first place is reachable depends mainly on the cap. Under a fixed gain per task the order of rules differs between the datasets and changes with the margin by which first place must be won (appendix), and on the Open LLM Leaderboard the order of the median and the minimum rule's cost in points also changes with the data slice (appendix; slice medians use different target sets). The median $h^\ast$ is reported by rank bin: for ranks 11--30 it is \hstarBinHelmMeanElevenThirty{} (HELM) and \hstarBinOllbMeanElevenThirty{} (Open LLM Leaderboard) points under the mean, against a pooled 95th-percentile fine-tuning gain of \empPooled{} points, while the all-target median is dominated by the \ollbShareOverHundred\% of Open LLM Leaderboard targets ranked beyond 100. Under per-task calibrated caps most targets cannot reach first place, and for HELM this depends on the transfer of the fine-tuning gains (\empMmluQ{} points from one task, \empPooled{} pooled; Section~\ref{sec:exp-caps}). We present this as a sensitivity, not as a prediction.
\item \textbf{Ordinal rules differ in kind.} On HELM (ranks 11--30, ceiling) plurality is cheapest in both measures, while Borda and the Condorcet requirement need many tasks; contrary to our expectation, the Condorcet requirement is cheaper than the mean in points (Section~\ref{sec:exp-ordinal}). On the Open LLM Leaderboard, for all targets of ranks 2--100, plurality is again the hardest rule to make feasible under a fixed gain (\ordHundredTenPluralityFeas\% of targets at cap $+10$).
\end{itemize}
We discuss the assumptions behind these results, what breaks them, and what we did not do in Section~\ref{sec:limits}. Related work is in Section~\ref{sec:related}.
